\subsection{代数闭域}

命题1. --- 设K为一个域；则以下性质等价：

(AC) \(\mathbf{K}[\mathbf{X}]\) 中的每个非常数多项式在 \(\mathbf{K}[\mathbf{X}]\) 中分裂成若干个一次多项式（可以相同也可以不同）的乘积。

(AC') K{[}X{]} 的每一个非常数多项式在K中至少有一个根。

(AC'') K{[}X{]} 中的每一个不可约多项式的次数都为1。

(AC''') K 的每个代数扩张的次数均为 1（换言之，K 在 K 的每个扩域中都是相对代数闭的）。

我们首先证明性质 (AC)、(AC') 与 (AC'') 的等价性。显然 (AC) 蕴含 (AC'')。由于 K{[}X{]} 的每个非常数多项式可被一个不可约多项式整除（IV，第13页，命题13），且 K{[}X{]} 中每个次数为 1 的多项式显然在K中有一个根，我们看到(AC'')蕴含(AC')。条件(AC')通过对n的归纳，蕴含 K{[}X{]} 中每个n次多项式都是n个一次多项式的乘积（IV，第14页，命题1），因此(AC')蕴含(AC)。

还需证明 \((AC'')\) 与 \((AC''')\) 是等价的。若 \((AC'')\) 成立，则K的扩域L中每个在K上代数的元素次数都为1（V，第16页，定理1），因此属于K，这便确立了 \((AC''')\) 。反之，设 f 是一个次数为 \(n \geqslant 1\) 的不可约多项式，它属于 K{[}X{]}；商代数 \(K[X]/(f)\) 在 K 上的次数为 n 且是一个域，因此是 K 的一个 n 次代数扩张（V，第 18 页，命题 2）。现在显然 \((AC''')\) 蕴含 \((AC'')\) 。

定义 1. --- 一个域 K 被称为代数闭的，如果它满足（等价的）性质 (AC)、(AC')、(AC'')、(AC''')。

* 例 1. --- 复数域 \(\mathbf{C}\) 是代数闭域（Gen.～Top., VIII, p.～100）。*

一个域 K 在扩域 E 中相对代数闭，并不一定意味着 K 是代数闭的（事实上，每个域在自身中都是相对代数闭的，并且存在不是代数闭的域，例如 Q 或 \(F_{p}\) * 或 \(R_{*}\) ）。然而：

命题2。——设 \(\Omega\) 为一个代数闭域，K 为 \(\Omega\) 的一个子域。那么 K 的相对代数闭包 \(\bar{\mathbf{K}}\) （在 \(\Omega\) 中）是一个代数闭域。

设 \(f\) 为 \(\bar{\mathbf{K}}[\mathbf{X}] \subset \Omega[\mathbf{X}]\) 中的一个非常数多项式。由于 \(\Omega\) 是代数闭的，多项式 \(f\) 在 \(\Omega\) 中至少有一个根，且由于这个根在 \(\bar{\mathbf{K}}\) 上是代数的，它属于 \(\bar{\mathbf{K}}\) （V，第19页，命题4）。因此 \(\bar{\mathbf{K}}\) 满足 (AC')。

例2.——根据命题2，所有在Q上代数的复数（通常简称为代数数）构成一个代数闭域。

每个代数闭域都是无限的。

设 \(\mathbf{K}\) 为一个有限域，并令 \(f(\mathbf{X}) = 1 + \prod_{a\in \mathbf{K}}(\mathbf{X} - a)\) 。多项式

\(f \in K[X]\) 是非常数的，并且 \(f(a) = 1\) 对每个 \(a \in K\) 成立。因此，域K不满足(AC')，从而不是代数闭的。

定理1（施泰尼茨）。——设K为一个域，E为K的一个代数扩张，且 \(\Omega\) 为K的一个代数闭扩张；则存在一个 E 到 \(\Omega\) 的 K-同态。

由V，第13页，附注，存在一个扩域 \(\Omega'\) （\(\Omega\) 的扩张）以及一个K-同态 \(u\) ，它将 E 映入 \(\Omega'\) 。设 \(x \in E\) ；由于 \(x\) 在K上是代数的， \(u(x)\) 在 \(u(\mathbf{K})\) 上是代数的，并且更有理由在 \(\Omega\) 上是代数的（V，第17页，推论2）；由于 \(\Omega\) 是代数闭的，因此我们有 \(u(x) \in \Omega\) 。由此可知， \(u\) 将 E 映入 \(\Omega\) 。

